\section{Metodología}
\label{sec:metodologia}

Este capítulo describe el conjunto de datos, el protocolo experimental, los parámetros del sistema y las métricas de evaluación empleadas para caracterizar el sistema MAE-TMS.

% ---------------------------------------------------------------
\subsection{Conjunto de datos}
\label{sec:dataset}

Se utilizó el conjunto de datos ETH-80~\cite{leibe2003eth80}, que contiene imágenes fotográficas de objetos tridimensionales.
Para este experimento se seleccionaron tres clases: \texttt{apple} (manzana), \texttt{horse} (caballo) y \texttt{car} (automóvil).
Cada clase cuenta con 410 imágenes en escala de grises de $128\times128$ píxeles.
Las imágenes de cada clase se dividieron en 328 imágenes de entrenamiento y 82 de prueba, sin solapamiento.

Dentro del conjunto de entrenamiento se definieron tres piscinas disjuntas con propósitos distintos:
\begin{itemize}
    \item \textbf{Llenado} (\textit{filling}): imágenes \texttt{train[0:200]} de cada clase, utilizadas para registrar patrones en las memorias homoasociativas de cada agente especialista.
    \item \textbf{Interacciones visuales}: imágenes \texttt{train[200:328]} de cada clase, usadas para simular las interacciones de percepción durante la fase de formación del directorio transactivo.
    \item \textbf{Evaluación}: las 82 imágenes de prueba (\texttt{test}) de cada clase, nunca vistas durante el entrenamiento ni la formación del directorio.
\end{itemize}

% ---------------------------------------------------------------
\subsection{Pipeline semántico}
\label{sec:pipeline-semantico}

Cada consulta en lenguaje natural transita por un pipeline de cuatro etapas antes de llegar a las memorias asociativas.

\paragraph{Extracción de conceptos con ConceptNet.}
Se consultó ConceptNet~5.7~\cite{speer2017conceptnet} con las relaciones \texttt{IsA}, \texttt{HasProperty}, \texttt{RelatedTo}, \texttt{HasPart}, \texttt{HasA}, \texttt{MadeOf} y \texttt{CapableOf}.
Solo se retuvieron aristas con peso mayor o igual a 2.0, lo que produjo un vocabulario de aproximadamente 20 etiquetas (\textit{labels}) por clase con su frecuencia de co-ocurrencia.

\paragraph{Vectorización con fastText.}
Cada etiqueta del vocabulario se representó con un vector de 300 dimensiones extraído del modelo fastText preentrenado~\cite{bojanowski2017fasttext}.
Para obtener un patrón binario compatible con la Memoria Asociativa Entrópica, cada vector se binarizó tomando el signo de cada componente, resultando en un vector de 300 bits por etiqueta.

\paragraph{Procesamiento de consultas con spaCy.}
Las consultas en lenguaje natural se preprocesaron con spaCy~\cite{honnibal2020spacy}: se aplicó lematización y se retuvieron únicamente tokens de categoría gramatical \texttt{NOUN}, \texttt{ADJ} o \texttt{PROPN}.
Se eliminaron palabras vacías (\textit{stopwords}) y se aplicó deduplicación estable para conservar el orden del original.

\paragraph{Normalización por lemas.}
Las etiquetas de ConceptNet se almacenan en su forma superficial (e.g., \texttt{wheels}, \texttt{hooves}), mientras que spaCy devuelve lemas (\texttt{wheel}, \texttt{hoof}).
Para evitar rechazos espurios, se construyó un índice de alias lema$\to$etiqueta de vocabulario que permite resolver la consulta aun cuando la forma superficial difiera del lema.

% ---------------------------------------------------------------
\subsection{Percepción visual: encoder ResNet-18}
\label{sec:percepcion}

Para obtener representaciones latentes de las imágenes se entrenó un autoencoder convolucional con encoder basado en ResNet-18~\cite{he2016resnet}.
El espacio latente tiene dimensión 64.
Los vectores latentes se cuantizaron a valores binarios mediante normalización global min-max con un margen del 10\%, calculada exclusivamente sobre el pool de llenado (\texttt{train[0:200]} por clase).
Esta cuantización global evita que la escala relativa entre clases afecte la codificación binaria.

% ---------------------------------------------------------------
\subsection{Banco de evaluación semántica}
\label{sec:banco}

Se construyó un banco de 80 consultas en lenguaje natural con etiqueta de clase correcta conocida (\textit{ground truth}): 27 consultas de \texttt{apple}, 27 de \texttt{horse} y 26 de \texttt{car}, presentadas de forma intercalada para reflejar un escenario de uso realista.
Las consultas varían en longitud y vocabulario (e.g., ``a red crunchy fruit'', ``an animal that trots'', ``a wheeled motor vehicle'') y cubren tanto descripciones directas como perifrásticas.

% ---------------------------------------------------------------
\subsection{Parámetros de la EAM}
\label{sec:parametros}

Los parámetros propios del modelo de Memoria Asociativa Entrópica de Pineda y Morales~\cite{pineda2021entropic,pineda2024hetero} se fijaron en:
$\iota = 0$, $\kappa = 0$, $\xi = 0$, $\sigma = 0.1$.
La justificación experimental de estos valores se presenta en la Sección~\ref{sec:param-iota-kappa} (barrido $\iota \times \kappa$) y en la Sección~\ref{sec:capacidad} (análisis de $\xi$).
En síntesis, $\iota > 0$ poda conocimiento escaso legítimo antes que promiscuidad (con patrones binarios las relaciones son ralas) y $\kappa$ resultó inerte a lo largo de todo el rango probado.

% ---------------------------------------------------------------
\subsection{Métricas de evaluación}
\label{sec:metricas}

Se emplearon las siguientes métricas:

\begin{itemize}
    \item \textbf{Precisión temprana} (\textit{early accuracy}): fracción de consultas para las que el agente ganador en la fase de activación temprana coincide con la clase real.
    \item \textbf{Precisión madura} (\textit{mature accuracy}): ídem tras la fase de recuperación madura, con lectura de directorio corregida por conteo (B1) o cruda.
    \item \textbf{Tasa de rechazo}: fracción de consultas para las que ningún agente supera el umbral de aceptación.
    \item \textbf{Fidelidad temprana-madura}: concordancia entre la clase ganadora en fase temprana y la clase ganadora en fase madura.
    \item \textbf{Entropía del directorio}: entropía de Shannon de la distribución de \textit{counts} de los agentes especialistas, normalizada respecto a $\log_2(3)$ (máximo para tres agentes).
    \item \textbf{Counts por agente}: número de registros acumulados en el directorio transactivo por cada agente especialista.
    \item \textbf{Routing de test} (hemisferio visual): fracción de imágenes de prueba asignadas correctamente por el directorio visual.
    \item \textbf{Evocación de dominio} (\textit{top-3 coseno}): fracción de imágenes para las que al menos una de las tres etiquetas más cercanas en el espacio fastText pertenece al vocabulario de la clase real (recall inverso imagen $\to$ palabras).
\end{itemize}

% ---------------------------------------------------------------
\subsection{Organización del experimento}
\label{sec:organizacion}

El experimento se organiza en cinco secciones (A--E) que caracterizan al sistema completo bajo un protocolo único, más un estudio de calibración del directorio que fundamenta la elección de lectura B1.
La Sección~\ref{sec:proto-completo} presenta el protocolo completo; la Sección~\ref{sec:param-iota-kappa} examina los parámetros nativos; la Sección~\ref{sec:formacion} analiza la formación del directorio; la Sección~\ref{sec:capacidad} estudia la curva de capacidad de llenado; y la Sección~\ref{sec:hemisferio-visual} evalúa el hemisferio visual.
